%% ==========================================================================
%%  FINAL REPORT --- SECTION 4, SYSTEM DESIGN
%%  Source: design/ARCHITECTURE-draft-v1.md sections 1-7 and its Decisions
%%  Log, corrected wherever the implementation settled a question differently.
%% ==========================================================================
\section{System design}
\label{sec:design}

\sys{} has three stages that form a cycle (Figure~\ref{fig:system}). Training
runs periodically over all users and produces the model. Serving and delivery
run per user, online. The harvest carries what users fetched back into the next
round of training.

\begin{figure}[t]
\centering
\begin{tikzpicture}[
    font=\small,
    box/.style={draw, rounded corners=2pt, align=center, inner sep=4pt,
                minimum height=10mm},
    note/.style={font=\scriptsize, align=left, text width=44mm},
    >={Stealth[length=5pt]}
  ]
  \node[box, minimum width=44mm] (train) at (0,0)
    {\textbf{Training} \quad $\party{0}, \party{1}, \party{2}$\\
     power iteration on $\shr{\ratings}$};
  \node[box, minimum width=44mm] (serve) at (7.2,0)
    {\textbf{Serving} \quad $\party{0}, \party{1}, \party{2}$\\
     $\shr{\mathbf{s}} = \shr{\avec{i}}\,\itemembed$, local};
  \node[box, minimum width=44mm] (user) at (7.2,-2.7)
    {\textbf{User}\\ reconstructs $\mathbf{s}$, picks top $\topk$};
  \node[box, minimum width=44mm] (pir) at (0,-2.7)
    {\textbf{Delivery} \quad $\party{0}, \party{1}$\\
     DPF-PIR, one key per fetch};

  \draw[->] (train) -- node[above, font=\scriptsize]{$\itemembed$ public,
    $\shr{\userembed}$ shared} (serve);
  \draw[->] (serve) -- node[right, font=\scriptsize]{score shares} (user);
  \draw[->] (user) -- node[above, font=\scriptsize]{$\topk$ DPF key pairs} (pir);
  \draw[->, dashed] (pir) -- node[left, font=\scriptsize, align=right]
    {harvest: shared\\ consumption counts} (train);

  \node[note, below=3mm of pir] {each server expands its key over all
    $2^{11}$ records, so the work never depends on the index};
  \node[note, below=3mm of user] {selection happens on the user's machine,
    so no server learns the ranking};
\end{tikzpicture}
\caption{The \sys{} cycle. Solid arrows run in every experiment in this report.
The dashed harvest arrow is implemented and tested as a library path
(Section~\ref{sec:design:harvest}) but is not yet driven by the networked server
binaries, and every training run here reads its ratings from disk.}
\label{fig:system}
\end{figure}

\subsection{Training}
\label{sec:design:train}

Training takes $\shr{\ratings}$ and runs \nudgeplain{}'s $\ApproxFactor$ under
replicated sharing: for each component, $\iters$ rounds of two shared
matrix--vector products, a projection against the rows of $\itemembed$ already
found, and a normalisation. Appendix~\ref{app:protocols} gives it as
Algorithm~\ref{alg:factor}. Three pieces of it need design decisions of their
own.

\paragraph{Where truncations go.} Ratings are integers, held at scale zero, so the
two matrix products need no truncation of their own, only a fixed rescaling
shift afterwards that keeps $\norm{\mathbf{v}}^2$ from overflowing. The remaining
truncations sit in orthogonalisation and normalisation, and each costs a
three-round protocol. Orthogonalisation can be \emph{deferred}: subtract the
projections at scale $3\fracbits$ and truncate once by $2\fracbits$. Or it can be
\emph{immediate}: truncate the inner products first and the result afterwards,
two truncations that never exceed scale $2\fracbits$. Deferral saves a
truncation and costs $\fracbits$ bits of headroom, and whether that fits depends
on $\bits$, $\fracbits$, $\items$ and $\dime$. So \sys{} derives the schedule from
public parameters at start-up, proves it fits with at least eight spare bits, and
refuses to run otherwise. At $\bits = 64$ on MovieLens-100K the deferred schedule
does not fit and training uses the immediate one, with 21 bits to spare. At
$\bits = 128$ the deferred schedule fits with 67.

\paragraph{How to normalise.} Normalising needs $1/\norm{\mathbf{v}}$ on a
shared $\mathbf{v}$. \sys{} has two implementations behind one interface. The
spec-faithful one follows \nudgeplain{}: a DCF-based gate finds the most
significant non-zero bit of $\norm{\mathbf{v}}^2$, which gives a first
approximation to $1/\norm{\mathbf{v}}$, and Newton--Raphson refines it. It
reveals nothing. The second opens $\norm{\mathbf{v}}^2$, divides in the clear,
and rescales by a public constant. It is far cheaper and it leaks, and
Section~\ref{sec:threat:gaps} states exactly what that costs the security
argument. Every result records which normaliser produced it.

\paragraph{Keeping $\itemembed$ public.} Because each finished row of
$\itemembed$ is opened, projecting against it is a product of shares with public
values, which needs no multiplication protocol. Our implementation priced the
alternative: keeping $\itemembed$ shared makes one small orthogonalisation cost
6 rounds and 1{,}800 bytes where the public version costs none. We keep
\nudgeplain{}'s choice, and Section~\ref{sec:leak} measures its consequence.

\subsection{Serving}
\label{sec:design:serve}

Each server holds a share of the user's embedding $\avec{i}$ and the public
$\itemembed$, so it computes its share of all $\items$ scores locally, with no
communication between servers. The user also supplies a shared mask carrying a
large negative sentinel at the items it has already rated. The servers add it to
their shares, which is again local, and send the result to the user, who
reconstructs the score vector and selects the top $\topk$.

The design draft specified an oblivious top-$\topk$ selector run by the servers.
We cut it: once the user reconstructs the scores, no server learns the selection
whether or not an oblivious selector exists, so the selector would protect
something already protected. The cost is that the user downloads a score share of
length $\items$ from each server, which Section~\ref{sec:limits} returns to.

\subsection{Delivery}
\label{sec:design:deliver}

Records are packed to a fixed width of $L = 256$ bytes, since a variable-width
record would leak its identity through the reply length. MovieLens-100K's 1{,}682
items pad to a domain of $2^{11} = 2{,}048$ entries, and padding records are all
zero, which leaks nothing because they are the same width as real ones. The DPF
is a two-party primitive, so $\party{0}$ and $\party{1}$ answer PIR queries and
$\party{2}$ does not. All three servers hold the catalogue, which is public. For
each of its $\topk$ recommendations the user sends one key to each PIR server,
and the difference of the two answers is the record.

\subsection{Closing the loop}
\label{sec:design:harvest}

When a PIR server expands a key, the expansion it computes is its share of the
indicator vector of the fetched item. \pirsonaplain{}'s observation is that this
is exactly a secret-shared record of consumption, already computed and normally
thrown away. \sys{}'s answer routine returns the expansion alongside the answer,
an accumulator sums it over queries, and a one-round conversion turns the two PIR
servers' 2-of-2 shares into the 2-of-3 replicated form training consumes. A
further one-round check, added after Section~\ref{sec:leak:malicious} found the
need for it, confirms that each query carries weight exactly one.
Algorithm~\ref{alg:harvest} gives both.

We state the extent of this plainly. The harvest, the conversion and the weight
check are implemented, and a test drives real PIR queries through them and
multiplies the result under the training substrate. The networked server
binaries do not yet accumulate expansions across sessions or retrain from them,
and every training run in this report reads MovieLens ratings from disk.

\subsection{Design decisions}
\label{sec:design:decisions}

Table~\ref{tab:decisions} collects the decisions that shape everything else,
each with the reason recorded when it was taken.

\begin{table}[t]
\centering
\small
\caption{The design decisions that matter most, from the project's decisions
log.}
\label{tab:decisions}
\begin{tabular}{@{}p{5.0cm} p{9.6cm}@{}}
\toprule
Decision & Reason \\
\midrule
Power iteration, not gradient descent & Matrix--vector products cost
  communication in the output length only, so only truncation and normalisation
  are expensive. This is \nudgeplain{}'s core idea. \\
Three servers, 2-of-3 replicated sharing, semi-honest & Matches \nudgeplain{}
  exactly, one fewer non-colluding party than \pirsonaplain{}'s four. \\
$\itemembed$ opened row by row & Makes orthogonalisation local. The cost is the
  leakage measured in Section~\ref{sec:leak}. \\
One DPF implementation for both halves & The same primitive is the PIR query and,
  through the DCF built beside it, the comparison gate. Forking it would double
  the code most likely to be wrong. \\
Top-$\topk$ on the user's machine & The user reconstructs the scores anyway, so a
  server-side oblivious selector would protect nothing further. \\
Fixed-width records & A variable width leaks the record through the reply
  length. \\
Ring width a template parameter & Whether $\bits = 64$ suffices was an open
  question, and the answer turned out to depend on which normaliser runs
  (Section~\ref{sec:eval:ring}). \\
Three OS processes, no threads & Each server serves one client to completion, so
  nothing is concurrent, and separate processes make that true rather than
  likely. \\
\bottomrule
\end{tabular}
\end{table}
